\section{Conclusion}
\label{sec:conclusion}

We characterized the post-training threshold-calibration stage in federated IoT anomaly detection as an attack surface that training-phase and aggregation-phase defenses structurally do not address.
Under a score-level contamination model with a single compromised client, the raising attack degrades victim detection rates by 2.4--7.9\,pp across non-zero injection fractions (4.9--7.9\,pp at $f{=}0.4$) with 100\,\% Gate-1 pass rates and 95\,\% bootstrap CIs excluding zero.
The lowering attack shifts detection thresholds downward, increasing victim TPR by 10--12\,pp and $\Delta\mathrm{CV(FPR)}$ by 0.05--0.06 (as a lower bound under \textsc{Replace-Fixed-Budget}) for \textsc{Local} and \textsc{Cluster}; both effects are operationally meaningful in IoT deployment contexts.
The policy ordering---per-victim exposure maximized under \textsc{Local}, diluted but federation-wide under \textsc{Global}, intra-cluster under \textsc{Cluster}---implies that calibration-stage hardening must match the aggregation design, not only the training-phase defense.
Future work should address traffic-level realizability of the score-level injection model, calibration-stage integrity defenses, and generalization across anomaly-detector architectures and IoT datasets.
